\section{模型评价、改进与推广}

\subsection{模型评价}

本文围绕微网购电与储能调度建立统一优化框架，以混合整数线性规划描述电量平衡、储电量变化及充放电互斥等约束，并逐步引入负载与光伏发电量预测误差、紧急购电、日内预报更新和波动电价。各问题均按照决策时点的信息边界进行求解，并通过储电量跨日传递保持运行状态连续，从而实现不同情形下购电与储能策略的统一比较。

模型同时考虑正常购电、购电调整和紧急购电费用，并设置完全信息基准，用于衡量预测误差带来的额外成本。通过不同预测方法、安全冗余和滚动调整策略的比较，可进一步分析信息更新、储能调度和电价变化对购电费用及紧急购电风险的影响。

模型仍存在一定局限。现有预测与策略比较主要基于2025年数据，其在其他年份和不同电价波动条件下的稳定性仍需进一步检验；问题三和问题四采用限制减少购电并保持日初储能方案不变的调整规则，该设定具有一定保守性，其适用范围与具体供需及价格特征有关。

\subsection{模型改进}

当前模型主要采用历史数据形成负载、光伏发电量和电价的点预测，对极端偏差及预测误差分布考虑较少。后续可利用历史残差构造不确定区间或典型场景，并引入随机优化或鲁棒优化，在控制购电成本的同时提高策略对预测偏差的适应能力。

现有模型主要考虑储能容量、充放电功率及效率约束，尚未计入长期循环使用产生的退化成本。进一步可根据充放电深度和循环次数构造储能寿命成本，并将其纳入目标函数，以兼顾短期购电费用和长期设备使用成本。

\subsection{模型推广}

更换负载、光伏发电量、储能参数及电价数据后，本文模型可应用于工业园区、商业建筑群等具有分布式发电、储能和外部电网购电需求的用户侧能源系统。若增加风电、可控负荷或电动汽车等资源，可在现有电量平衡和运行约束基础上扩展相应变量，实现多类型能源资源的联合调度。

进一步考虑多个微网之间的电能交易和需求响应后，可将单一微网购电问题扩展至区域多微网协同调度，为分布式能源消纳、储能运行和市场化购电提供决策支持。